%%%%%%%%%%%%%%%%%%%%%%%%%%%%%%%%%%%%%%%%%%%%%%%%%%%%%%%%%%%%%%%%%%%%%%%%%%%%%%%%
% chiral_paper.tex
% A chiral quantum-state hypothesis for cosmic expansion and spiral galaxies
% A4 format, numeric citations [n], Nature-style manuscript
%%%%%%%%%%%%%%%%%%%%%%%%%%%%%%%%%%%%%%%%%%%%%%%%%%%%%%%%%%%%%%%%%%%%%%%%%%%%%%%%

\documentclass[a4paper,10pt]{article}

% Page geometry
\usepackage[utf8]{inputenc}
\usepackage[T1]{fontenc}
\usepackage{geometry}
\geometry{a4paper,left=20mm,right=20mm,top=25mm,bottom=25mm}

% Maths and symbols
\usepackage{amsmath,amsfonts,amssymb}

% Figures
\usepackage{graphicx}
\usepackage{hyperref}

% Line numbers for manuscript style
\usepackage{lineno}
\linenumbers

% Author/affiliation formatting
\usepackage{authblk}
\renewcommand\Authands{ and }

% Custom keywords command
\providecommand{\keywords}[1]{%
  \vspace{0.5em}\noindent\textbf{Keywords:} #1\par\vspace{0.5em}}

\title{\textbf{A chiral quantum-state hypothesis for cosmic expansion and the
       spiral morphology of galaxies}}

\author{Xiaohui Li}
\affil{Independent Researcher, Baoji 721000, China}
\affil{Correspondence: \texttt{danceluck@163.com};
       ORCID: 0009-0003-8233-7191}

\date{}

\begin{document}

\maketitle

\begin{abstract}
The Big Bang is conventionally treated as an initial condition, and the origin
of angular momentum in spiral galaxies remains an open problem in
astrophysics. We propose a unified, falsifiable conjecture: fundamental
quantum states carry a left--right chirality, and identical chiral states
attract whereas opposite chiral states repel. In the earliest universe, an
equal mixture of left- and right-handed states produces a net repulsive
interaction that we identify, phenomenologically, with the expanding phase of
the Big Bang. A minute stochastic fluctuation favouring one handedness then
triggers a chiral phase transition; the resulting single-chiral background
mediates an effective long-range attraction interpretable as gravitation,
while its helical topology seeds the vorticity that produces rotating, spiral
galaxies. We formulate a minimal field-theoretic model, derive six testable
signatures, and outline a validation programme using Planck, SDSS, Galaxy
Zoo, Gaia and IllustrisTNG data. The framework is presented as a theoretical
hypothesis awaiting empirical confirmation.
\end{abstract}

\keywords{chirality; quantum states; cosmic expansion; spiral galaxies;
large-scale structure; symmetry breaking}

\section{Introduction}

The standard $\Lambda$CDM cosmology successfully describes cosmic expansion,
the cosmic microwave background (CMB), and the formation of large-scale
structure \cite{planck18,peebles,dodelson}. Yet it leaves two deep questions
unanswered. First, why did the universe begin in a hot, dense, expanding
state---the Big Bang? Inflation provides a mechanism for the origin of
structure \cite{linde}, but the initial expansion itself is usually imposed by
hand. Second, why do most galaxies exhibit ordered rotation and spiral arms?
Tidal torque theory explains the acquisition of angular momentum from
large-scale tides \cite{leepen,codis}, but the statistical distribution of
spin directions and any possible cosmological handedness remain debated
\cite{longo,iye}.

Here we propose a speculative but conceptually simple unification. We
postulate that quantum states possess an intrinsic left--right chirality and
that the interaction between two chiral states follows a
``same-attract, opposite-repel'' rule. The idea is motivated by the behaviour
of magnetic field lines: filings lying on the same field line are mutually
attracted into a filament, whereas field lines of the same sense repel each
other. We translate this magnetic analogy into a ``chiral matter-line''
model.

If the primordial universe contained equal numbers of left- and right-handed
states, the dominant inter-chiral repulsion would drive an outward
expansion---identified here with the Big Bang. A tiny fluctuation that made
one chirality slightly more numerous would be amplified by a chiral phase
transition, leaving the universe dominated by a single handedness. The
surviving same-chiral background then provides an attractive force analogous
to gravity, and its helical structure naturally seeds the rotation and spiral
morphology of galaxies. We emphasize that this is a hypothesis, not an
established theory; we supply concrete tests by which it can be falsified.

\section{Results}

\subsection{A minimal chiral matter-line model}

Consider a continuum of quantum ``matter-line'' elements labelled by chirality
$a\in\{L,R\}$. The pairwise interaction energy between two elements separated
by distance $r$ is
\begin{equation}
U_{ab}(r)=-\alpha\,\delta_{ab}\frac{e^{-r/\lambda}}{r}
+\beta\,(1-\delta_{ab})\frac{e^{-r/\lambda}}{r},
\label{eq:Uab}
\end{equation}
where $\alpha>0$ governs same-chiral attraction, $\beta>0$ governs
opposite-chiral repulsion, and $\lambda$ is a coherence length. The Kronecker
delta $\delta_{ab}$ selects the interaction channel. Equation~\eqref{eq:Uab}
is the simplest Yukawa-like form consistent with the postulated rule.

In a homogeneous early universe with number densities $n_L$ and $n_R$, the
interaction energy density is
\begin{equation}
\mathcal{E}_{\rm int}
=\frac12\sum_{a,b}n_a n_b\int d^3r\,U_{ab}(r)
\approx -\frac14\alpha\lambda^2(n_L^2+n_R^2)
+\frac12\beta\lambda^2 n_L n_R.
\label{eq:Eint}
\end{equation}
For an exactly symmetric mixture, $n_L=n_R=n/2$, this becomes
\begin{equation}
\mathcal{E}_{\rm int}=\frac{n^2\lambda^2}{8}(\beta-\alpha).
\label{eq:Esym}
\end{equation}
If $\beta>\alpha$, the total interaction energy is positive. The system can
lower its energy by separating the two chiral species and increasing the
volume they occupy. We interpret this outward segregation as the kinetic
driver of the initial cosmic expansion---the Big Bang in this picture is the
macroscopic manifestation of primordial chiral repulsion.

\subsection{Chiral symmetry breaking}

Define the order parameter $\eta=(n_L-n_R)/(n_L+n_R)$. Near a critical
temperature $T_c$ the free-energy density can be written as
\begin{equation}
f=f_0+a(T-T_c)\eta^2+b\eta^4,\qquad a,b>0.
\label{eq:GL}
\end{equation}
For $T<T_c$ the symmetric state $\eta=0$ becomes unstable and the minima occur
at $\eta=\pm1$. Which minimum is selected depends on the first statistically
significant fluctuation; therefore the final handedness of the universe is
random. This naturally explains why the observed cosmos has a single dominant
chirality without requiring a pre-existing bias: a minute excess, possibly
quantum in origin, is amplified to global symmetry breaking.

\subsection{Emergent gravitation in a single-chiral background}

After symmetry breaking, only one chiral species remains as the dominant
background. A test mass $m_i$ embedded in this background interacts
attractively with another mass $m_j$ through the same-chiral channel of
Eq.~\eqref{eq:Uab}. Averaging over the uniform background density $n_{\rm bg}$
gives an effective two-body potential
\begin{equation}
V_{\rm eff}(r)=
-\frac{\alpha n_{\rm bg}\lambda^2}{m^2}\,
\frac{m_i m_j}{r}\,e^{-r/\lambda}.
\label{eq:Veff}
\end{equation}
If $\lambda$ is much larger than cosmological scales, the Yukawa cutoff can be
ignored. Identifying the prefactor with Newton's constant yields the
consistency relation
\begin{equation}
G \approx \frac{\alpha n_{\rm bg}\lambda^2}{m^2}.
\label{eq:G}
\end{equation}
Equation~\eqref{eq:G} does not uniquely fix $\alpha$, $n_{\rm bg}$ and
$\lambda$, but it shows that the observed weakness of gravity ($G$ small) can
arise if the chiral background density is low or the coherence length is
large. The model therefore predicts that gravitational attraction is a
collective, emergent property of a single-chiral quantum vacuum.

\subsection{Spiral structure and angular momentum}

A uniform chiral background is not rotationally trivial: it carries helicity
density
\begin{equation}
\mathcal{H}=\mathbf{A}\cdot(\nabla\times\mathbf{A}),
\label{eq:helicity}
\end{equation}
where $\mathbf{A}$ is a vector potential whose circulation defines the
handedness of the matter lines. A left-handed background has $\mathcal{H}<0$
and a right-handed background has $\mathcal{H}>0$. When a density perturbation
collapses in this helical medium, the chiral field exerts a torque
\begin{equation}
\boldsymbol{\tau}\approx\int\rho\,\mathbf{r}\times
\mathbf{g}_{\rm chiral}\,d^3r,
\qquad
\mathbf{g}_{\rm chiral}\propto\mathcal{H}\,\mathbf{r},
\label{eq:torque}
\end{equation}
so that the resulting halo acquires angular momentum $\mathbf{J}$ whose sign
is locked to the background helicity. Conservation of angular momentum during
collapse then produces rotating disks. Spiral arms are density waves
\cite{linshu} superposed on the rotating disk; their handedness is correlated
with the sign of $\mathcal{H}$. Thus the prevalence of spiral galaxies
reflects the single chirality of the cosmic background.

\subsection{Testable predictions}

The model makes the following falsifiable predictions:
\begin{enumerate}
\item \textbf{Large-scale parity asymmetry.} The distribution of spiral-galaxy
  handedness should exhibit a dipole anisotropy, with amplitude
  $10^{-3}$--$10^{-2}$, detectable in Galaxy Zoo \cite{galaxyzoo},
  Iye--Sugai \cite{iye} and Gaia DR3 \cite{gaia} samples.
\item \textbf{Spin--tide correlation.} The spin axes of haloes should correlate
  more strongly with the local cosmic-web tidal shear than predicted by
  standard tidal-torque theory \cite{leepen,codis}. The excess signal scales
  with $\lambda$ and can be measured in N-body simulations.
\item \textbf{CMB parity anomalies.} A helical vacuum imprints parity-odd
  fluctuations on the CMB. The model predicts enhanced low-$\ell$ power
  asymmetry and non-zero $\langle T(\mathbf{n})E(\mathbf{n}')B(\mathbf{n}'')
  \rangle$ correlations, testable with Planck data \cite{planck18}.
\item \textbf{Chiral primordial gravitational waves.} The same helicity
  density that seeds vorticity should source circularly polarized gravitational
  waves, producing a non-zero TB or EB cross-correlation in CMB B-modes
  \cite{bicep}.
\item \textbf{Cosmic vorticity.} The dark-matter velocity field should possess
  a coherent curl component at the level $|\nabla\times\mathbf{v}|\sim
  10^{-14}$ s$^{-1}$, accessible through redshift-space distortions
  \cite{sdss}.
\item \textbf{Lepton chiral asymmetry.} The phase transition that fixes cosmic
  chirality may leave a relic neutrino--antineutrino asymmetry of order
  $10^{-9}$--$10^{-10}$, influencing big-bang nucleosynthesis and the CMB
  effective number of relativistic species \cite{goldhaber,christenson}.
\end{enumerate}

\section{Methods}

\subsection{Datasets}

We propose to test the hypothesis against publicly available data:
Planck 2018 CMB temperature and polarization maps and angular power spectra
\cite{planck18}; SDSS-IV eBOSS DR16 galaxy clustering and redshift-space
distortion measurements \cite{sdss}; Galaxy Zoo 2 morphological
classifications of spiral handedness \cite{galaxyzoo}; Gaia DR3 astrometry and
radial velocities for local spiral-chirality mapping \cite{gaia}; and
IllustrisTNG cosmological hydrodynamical simulations for benchmark halo spin
statistics \cite{illustris,illustris2}.

\subsection{Simulation strategy}

The chiral interaction can be implemented in the N-body code GADGET-4
\cite{gadget4} by adding a long-range Yukawa term to the gravitational force,
with sign controlled by a particle chirality tag. We recommend simulation
suites with box sizes $L=100$, $300$ and $1000\,h^{-1}$ Mpc and $N=1024^3$
particles, run for both $\Lambda$CDM and the chiral model. Initial conditions
should be generated with MUSIC \cite{music} using the same random seed.
Benchmarks:
(i) matter power spectrum:
$|P_{\rm chiral}(k)-P_{\Lambda\rm CDM}(k)|/P_{\Lambda\rm CDM}(k)<2\%$ for
$k<1\,h$ Mpc$^{-1}$;
(ii) halo mass function: agreement within $10\%$ for
$M>10^{12}\,M_\odot$;
(iii) galaxy spin dipole amplitude: $\Lambda$CDM predicts $A<2\times10^{-2}$;
the chiral model predicts $A\sim10^{-2}$--$10^{-3}$.

\subsection{Statistical validation}

Spiral handedness will be binned in a HEALPix map. The dipole amplitude and
direction will be extracted by a $\chi^2$ fit to
$P(\theta,\phi)=P_0[1+A\cos\theta']$. Significance will be assessed by
randomizing handedness $10^6$ times and computing a $p$-value. The same
procedure will be applied to simulated galaxy catalogs.

\section{Data availability}

All observational datasets used in the proposed validation are publicly
available through the Planck Legacy Archive, SDSS Science Archive Server,
Galaxy Zoo data repository, Gaia Archive, and IllustrisTNG database.
Simulation outputs and derived catalogs will be released upon acceptance.

\section{Code availability}

The modified GADGET-4 patch, initial-condition generation scripts, and
analysis pipelines will be deposited at Zenodo and GitHub
(\url{https://github.com/xiaohui-li-lab/}) with version-controlled
DOIs.
Please check the author's papers from 10-27 to 11-5 recently（2026-10-27 to 2026-11-5 in Github）; this forms a series.


\section{Author contributions}

X.L. conceived the chiral matter-line hypothesis, developed the formal model,
designed the validation programme, wrote the manuscript, prepared all figures,
and approved the final version.

\section{Competing interests}

The author declares no competing interests.

\section{Acknowledgements}

The author thanks the open-science community for making Planck, SDSS, Galaxy
Zoo, Gaia and IllustrisTNG data publicly accessible. Discussions with
colleagues in cosmology and quantum foundations helped shape the presentation
of this hypothesis.

\begin{thebibliography}{20}

\bibitem{planck18}
Planck Collaboration,
``Planck 2018 results. VI. Cosmological parameters'',
\textit{Astron. Astrophys.} \textbf{641}, A6 (2020).

\bibitem{sdss}
Alam, S. et al. (SDSS-IV),
``The completed extended Baryon Oscillation Spectroscopic Survey'',
\textit{Astrophys. J. Suppl.} \textbf{259}, 35 (2022).

\bibitem{galaxyzoo}
Lintott, C. J. et al.,
``Galaxy Zoo: morphologies derived from visual inspection of galaxies from the
Sloan Digital Sky Survey'',
\textit{Mon. Not. R. Astron. Soc.} \textbf{389}, 1179--1189 (2008).

\bibitem{gaia}
Gaia Collaboration,
``Gaia Data Release 3. Summary'',
\textit{Astron. Astrophys.} \textbf{674}, A1 (2023).

\bibitem{illustris}
Nelson, D. et al.,
``The IllustrisTNG Simulations: public data release'',
\textit{Comput. Astrophys. Cosmol.} \textbf{6}, 2 (2019).

\bibitem{illustris2}
Springel, V. et al.,
``First results from the IllustrisTNG simulations: matter and galaxy
clustering'',
\textit{Mon. Not. R. Astron. Soc.} \textbf{475}, 676--698 (2018).

\bibitem{longo}
Longo, M. J.,
``Evidence for a dipole in the spin axes of spiral galaxies'',
\textit{Astrophys. J. Lett.} \textbf{741}, L19 (2011).

\bibitem{iye}
Iye, M. \& Sugai, H.,
``A search for dipole anisotropy in the handedness of spiral galaxies'',
\textit{Publ. Astron. Soc. Jpn} \textbf{63}, L43--L46 (2011).

\bibitem{peebles}
Peebles, P. J. E.,
\textit{Principles of Physical Cosmology}
(Princeton Univ. Press, 1993).

\bibitem{dodelson}
Dodelson, S. \& Schmidt, F.,
\textit{Modern Cosmology}, 2nd edn
(Academic Press, 2021).

\bibitem{ryden}
Ryden, B.,
\textit{Introduction to Cosmology}, 2nd edn
(Cambridge Univ. Press, 2017).

\bibitem{goldhaber}
Goldhaber, M., Grodzins, L. \& Sunyar, A.,
``Helicity of neutrinos'',
\textit{Phys. Rev.} \textbf{109}, 1015--1017 (1958).

\bibitem{christenson}
Christenson, J. H. et al.,
``Evidence for the $2\pi$ decay of the $K_2^0$ meson'',
\textit{Phys. Rev. Lett.} \textbf{13}, 138--140 (1964).

\bibitem{linde}
Linde, A. D.,
``Inflationary cosmology'',
\textit{Lect. Notes Phys.} \textbf{738}, 1--54 (2008).

\bibitem{bicep}
Ade, P. A. R. et al. (BICEP/Keck),
``Improved constraints on primordial gravitational waves using Planck, WMAP,
and BICEP/Keck observations through the 2018 observing season'',
\textit{Phys. Rev. Lett.} \textbf{127}, 151301 (2021).

\bibitem{gadget4}
Springel, V. et al.,
``Simulating cosmic structure formation with the GADGET-4 code'',
\textit{Mon. Not. R. Astron. Soc.} \textbf{506}, 2871--2896 (2021).

\bibitem{leepen}
Lee, J. \& Pen, U.-L.,
``The spin of spiral galaxies and the cosmic web'',
\textit{Astrophys. J.} \textbf{681}, 1068--1074 (2008).

\bibitem{codis}
Codis, S. et al.,
``Connecting the spin of dark matter haloes to the cosmic web'',
\textit{Mon. Not. R. Astron. Soc.} \textbf{427}, 3320--3336 (2012).

\bibitem{linshu}
Lin, C. C. \& Shu, F. H.,
``On the spiral structure of disk galaxies'',
\textit{Astrophys. J.} \textbf{140}, 646--655 (1964).

\bibitem{music}
Hahn, O. \& Abel, T.,
``Multi-scale initial conditions for cosmological simulations'',
\textit{Mon. Not. R. Astron. Soc.} \textbf{415}, 2101--2121 (2011).

\end{thebibliography}
\textbf{\textbf{Declarations}}

  \textbf{\textbf{Acknowledgements}}

The author thanks colleagues for discussions on quantum aggregates, consciousness, emergence, and dimensionality. 

 \textbf{\textbf{Funding Statement}}:This research did not receive any specific grant from funding agencies in the public,commercial,or not-for-profit sectors.

Currently, there are no sponsors. Or I haven't found any sponsors or organizations yet. However, if I can publish it, I will cover part of the costs myself. Or it will be easier to find sponsors. (For example, the German company E.ON. I have previously found international buyers for their wind turbines.)

 \textbf{\textbf{Competing Interests}}

The author declares no competing interests.

 \textbf{\textbf{Author Contributions}}

Xiaohui Li: Conceptualization, Methodology, Investigation, Formal analysis, Writing – Original Draft, Visualization, Writing – Review \& Editing.

AI assisted me in rewriting the original text into an article that meets the submission requirements of Journal of the American Mathematical Society. No content was added or omitted.

 The author is solely responsible for the work described in this manuscript and has approved the final version for publication.

 \textbf{\textbf{Clinical trial:}}not applicable.

 \textbf{\textbf{Ethics and Consent to Participate declarations}}: not applicable.

 Hi,

This is an original work by the author. There is only one author:Xiaohui Li.

The article accurately reflects the scientific results.

Xiaohui Li agreed to submit the manuscript.

This manuscript has not been published in any other journals.

The author accepts the established procedure for selecting and publishing manuscripts.

good

Xiaohui Li



 \textbf{\textbf{Clinical trial number}}: not applicable.
\end{document}
